% Fig3 mechanism schematic — iter2 (post cross-figure consistency review)
% Key reframe: (b) shows DISCRETE (low-S, single/deep) -> CONTINUOUS (S80 broad) EVOLUTION,
% which is the paper's mechanism (CvS_Analysis_Strategy: discrete double-exp -> continuous power-law).
% n=breadth; rho60s separate (magnitude); carrier sign-agnostic; energy horizontal (DOS).
\documentclass[border=6pt]{standalone}
\usepackage{polymer-paper-preamble}
\usetikzlibrary{decorations.pathmorphing, arrows.meta}

\begin{document}
\begin{tikzpicture}[
  font=\sffamily\fontsize{7}{8.4}\selectfont,
  trapx/.style={cGray!62!black, font=\fontsize{5.5}{5.5}\selectfont},
  walk/.style={cGray!72!black, line width=1.05pt, rounded corners=1.3pt},
  carrier/.style={circle, fill=cGray!60!black, inner sep=0pt, minimum size=1.35mm},
  xfer/.style={-{Stealth[length=1.0mm]}, cGray!68!black, line width=0.6pt},
  lbl/.style={align=center, font=\fontsize{6.5}{8}\selectfont},
  sub/.style={align=center, font=\fontsize{6}{7.4}\selectfont},
  pl/.style={font=\fontsize{7.5}{9}\selectfont, font=\bfseries},
]
% ===================== (a) multiple-trapping ENERGY diagram: the carrier is repeatedly
%   captured into traps of varying depth and thermally re-emitted as it transits the
%   biased film -> dispersive transport, so I(t) decays. energy vertical, transit horizontal.
\node[pl] at (0.1,4.3) {a};
% the measured sulfur-polymer film (the sample) between the contacts -- names the material
\fill[cAmber!8] (0.72,1.02) rectangle (4.46,3.52);
% trap manifold: the continuous distribution that the discrete ticks below merely sample
\fill[cGray!10] (0.95,1.55) rectangle (4.28,2.82);
\node[anchor=north, font=\fontsize{5.5}{6.5}\selectfont, text=cAmber!55!black] at (2.59,0.98)
  {sulfur polymer film};
\draw[-{Stealth[length=1.0mm]}, line width=0.4pt] (0.42,0.95) -- (0.42,4.05);
\node[rotate=90, anchor=south, font=\fontsize{6.2}{7.5}\selectfont] at (0.20,2.35)
  {energy, $E$};
% biased film: contacts left (+V) / right (-V)
\fill[cGray!70] (0.56,1.05) rectangle (0.72,3.5);
\fill[cGray!70] (4.46,1.05) rectangle (4.62,3.5);
\node[font=\fontsize{6}{7.4}\selectfont] at (0.64,3.74) {$+V$};
\node[font=\fontsize{6}{7.4}\selectfont] at (4.54,3.74) {$-V$};
% transport level (mobility edge), tilted down by the applied field
\draw[cAmber!60!black, line width=0.8pt] (0.72,3.4) -- (4.46,3.0);
\node[anchor=west, font=\fontsize{5.5}{6.5}\selectfont, text=cAmber!55!black] at (1.0,3.62) {mobile transport level};
% trap states sampling the manifold: irregular depths/spacing = a continuous distribution
% (its spread = panel b's g(E)). The named states below are the only states that
% carrier paths may enter or leave; this prevents a visually detached arrowhead.
\coordinate (trap_shallow) at (1.55,2.45);
\coordinate (trap_mid) at (2.75,2.18);
\coordinate (trap_deep) at (3.60,1.75);
\foreach \x/\y in {1.55/2.45, 2.75/2.18, 3.60/1.75, 1.18/2.62, 1.95/1.98, 2.12/2.6,
                   2.45/2.0, 3.12/2.52, 3.32/2.05, 2.3/1.72, 1.42/2.08, 4.02/2.32}{
  \draw[cGray!55!black, line width=0.55pt] (\x-0.11,\y) -- (\x+0.11,\y);
}
\node[anchor=north, font=\fontsize{5.5}{6.5}\selectfont, text=cGray!55!black] at (2.5,1.46) {trap states};
% fig-agent:start object=carrier_walk panel=A kind=path truth_bearing=true
% Named carrier/trap endpoints bind every curved path to the state it depicts.
\node[carrier] (carrier_shallow_in) at (1.18,3.26) {};
\node[carrier] (carrier_shallow_out) at (1.88,3.10) {};
\node[carrier] (carrier_mid_in) at (2.34,3.12) {};
\node[carrier] (carrier_mid_out) at (3.06,2.92) {};
\node[carrier] (carrier_deep_in) at (3.44,2.88) {};
% shallow capture followed by thermal release
\draw[xfer] (carrier_shallow_in) to[out=-22,in=82] (trap_shallow);
\draw[xfer] (trap_shallow) to[out=82,in=-150] (carrier_shallow_out);
\node[font=\fontsize{4.8}{5.6}\selectfont, text=cGray!55!black, anchor=west] at (1.92,2.84) {release};
% mid-depth capture followed by thermal release
\draw[xfer] (carrier_mid_in) to[out=-20,in=84] (trap_mid);
\draw[xfer] (trap_mid) to[out=82,in=-152] (carrier_mid_out);
% deep capture; the carrier remains retained at this state in the schematic
\draw[xfer] (carrier_deep_in) to[out=-118,in=82] (trap_deep);
% fig-agent:end object=carrier_walk
\node[lbl, anchor=west, font=\fontsize{5.5}{6.8}\selectfont] at (0.45,0.55)
  {capture / release: current decays; resistance rises};
% fig-agent:start object=current_sparkline panel=A kind=inset_plot truth_bearing=true
\draw[-{Stealth[length=0.6mm]}, line width=0.3pt] (3.90,0.08) -- (4.34,0.08);
\draw[-{Stealth[length=0.6mm]}, line width=0.3pt] (3.90,0.08) -- (3.90,0.48);
\draw[cGray!72!black, line width=0.5pt] plot[smooth, domain=0:0.92, samples=32]
  ({3.90+0.48*\x},{0.08+0.34*exp(-3.2*\x)});
% fig-agent:end object=current_sparkline

% ===================== (b) g(E): DISCRETE -> CONTINUOUS evolution =====================
\node[pl] at (5.7,4.5) {b};
\begin{scope}[shift={(6.3,1.05)}]
  \draw[-{Stealth[length=1.1mm]}, line width=0.4pt] (0,0) -- (5.0,0) node[right, font=\fontsize{6.5}{8}\selectfont]{trap energy, $E$};
  \draw[-{Stealth[length=1.1mm]}, line width=0.4pt] (0,0) -- (0,3.2) node[rotate=90, anchor=south, xshift=-0.4mm, font=\fontsize{6.5}{8}\selectfont]{$g(E)$};
  % A schematic distribution: it compares discrete and continuous regimes, not band shapes.
  \draw[cBlue!65!black, line width=0.95pt] plot[smooth, domain=2.2:4.0, samples=60] (\x, {2.55*exp(-((\x-3.1)^2)/0.05)});
  % S80 = CONTINUOUS broad distribution, dashed red
  \draw[cRed!70!black, line width=0.95pt, dash pattern=on 2pt off 1.4pt] plot[smooth, domain=0.4:4.7, samples=90] (\x, {2.35*exp(-((\x-2.5)^2)/1.5)});
  % The physical layout places S60 on the right and S80 on the left, so the
  % composition cue must point right-to-left rather than imply S80 -> S60.
  \draw[-{Stealth[length=0.7mm]}, line width=0.3pt, cGray!45!black] (3.45,3.00) -- (1.55,3.00);
  \node[font=\fontsize{5.2}{6.3}\selectfont, text=cGray!55!black] at (2.50,3.16)
    {sulfur content $\uparrow$};
  % Distribution breadth follows the visible support of the broad distribution; it is
  % qualitative and is not directly equated with n.
  % fig-agent:start object=n_breadth panel=B kind=measurement_span truth_bearing=true
\draw[line width=0.35pt, cRed!60!black] (0.55,-0.25) -- (4.45,-0.25);
\draw[line width=0.35pt, cRed!60!black] (0.55,-0.12) -- (0.55,-0.38);
\draw[line width=0.35pt, cRed!60!black] (4.45,-0.12) -- (4.45,-0.38);
  \node[sub, cRed!60!black] at (2.5,-0.52) {distribution breadth};
  % fig-agent:end object=n_breadth
  % rho60s is an orthogonal vertical cue, deliberately not the height of either g(E) curve.
  \draw[{Stealth[length=0.7mm]}-{Stealth[length=0.7mm]}, line width=0.3pt, cGray!40!black] (4.55,0.10) -- (4.55,1.30);
  \node[font=\fontsize{5.0}{6.0}\selectfont, text=cGray!55!black, anchor=south, align=center]
    at (4.55,1.42) {$\rho_{60\mathrm{s}}$ magnitude};
  % Curve labels stay close to their visual referents; the colour and line style
  % are redundant cues, not the only encoding of discrete versus continuous traps.
  \node[font=\fontsize{5.3}{6.3}\selectfont, text=cBlue!60!black, align=left, anchor=west]
    at (3.48,2.38) {S60\\discrete trap};
  \node[font=\fontsize{5.3}{6.3}\selectfont, text=cRed!65!black, align=left, anchor=west]
    at (0.56,2.34) {S80\\continuous broad};
\end{scope}
\end{tikzpicture}
\end{document}
